% T4.6 — new Section III-B, "Scale factor and its compensation".
% Answers Reviewer 4 ("CORDIC gain/scaling compensation should be explained in
% details") and simultaneously carries contribution C1 against Reviewer 3.
%
% Place after the datapath/pipeline subsection and before the quantisation and
% synthesis methodology subsection.
%
% <<TBD>> marks numbers awaiting synthesis runs 1-4 in SYNTHESIS_RUN_LIST.md.

\subsection{Scale Factor and Its Compensation}
\label{sec:gain}

A shift-add micro-rotation is not a pure rotation. Stage $i$ of the Binary UD
datapath computes
\begin{equation}
  \begin{bmatrix} x_{i+1} \\ y_{i+1} \end{bmatrix}
  =
  \begin{bmatrix} 1 & -\delta_i 2^{-(i+1)} \\ \delta_i 2^{-(i+1)} & 1 \end{bmatrix}
  \begin{bmatrix} x_i \\ y_i \end{bmatrix},
  \label{eq:stage}
\end{equation}
which rotates by $\arctan(2^{-(i+1)})$ and simultaneously scales by
$\sqrt{1 + 2^{-2(i+1)}}$. Over $N$ stages the accumulated scale factor is
\begin{equation}
  K = \prod_{i=0}^{N-1}\sqrt{1 + 2^{-2(i+1)}},
  \label{eq:gain}
\end{equation}
which for $N = 10$ evaluates to $K = 1.164435$. Conventional CORDIC removes
this with a compensation multiplier or with additional scaling
micro-rotations, both of which cost area in the datapath.

\subsubsection{Why RoPE does not pay for it}
Because $\delta_i \in \{-1, +1\}$ at every stage --- the Binary UD formulation
admits no zero digit --- the factor in~\eqref{eq:gain} contains no dependence
on $\delta_i$, and $K$ is therefore the same for every rotation angle. This is
the property that makes compensation free in RoPE specifically. Queries and
keys are rotated by the same hardware, so a rotated query at position $m$ and
key at position $n$ are $K\mathbf{q}_m$ and $K\mathbf{k}_n$, and the attention
logit becomes
\begin{equation}
  \frac{\langle K\mathbf{q}_m,\, K\mathbf{k}_n \rangle}{\sqrt{d_h}}
  = K^2 \cdot \frac{\langle \mathbf{q}_m, \mathbf{k}_n \rangle}{\sqrt{d_h}} .
  \label{eq:fold}
\end{equation}
Since $K^2$ is a constant independent of $m$, $n$ and the feature index, it can
be absorbed into the existing $1/\sqrt{d_h}$ scaling that the attention block
already applies, by replacing it with $K^2/\sqrt{d_h}$ at no hardware cost. No
compensation stage is required anywhere in the rotator. We verify this in
simulation: across the full angle range the measured output magnitude is
$1.164426$ against the predicted $1.164435$, with a spread of $\pm 0.0003$ at
$W = 16$, confirming that the gain carries no angular dependence.

\subsubsection{Why stage merging forfeits the property}
The same argument fails for the merged radix-4 datapath of
Section~\ref{sec:merged}. A merged stage
applies $t = v \cdot 2^{-2i}$ with $v$ drawn from a digit set containing more
than one magnitude, so its gain is $\sqrt{1 + v^2 4^{-i}}$ and depends on which
digit was selected --- that is, on the angle. For the digit set of~\eqref{eq:merged},
$\{0, \pm 1, \pm \tfrac{3}{2}\}$, the accumulated gain ranges over $[1.00, 1.93]$. No
constant can absorb an angle-dependent factor, so~\eqref{eq:fold} no longer
applies and the distortion appears as a position- and dimension-dependent
rescaling of $\mathbf{q}$ and $\mathbf{k}$.

This is not an artefact of the particular digit set. A gain independent of the
digit requires every digit to share a magnitude, and a single-magnitude digit
set carries one bit per stage --- which is exactly the Binary UD datapath.
Constant gain and radix-4 stage merging cannot be obtained together.

\subsubsection{Recovering it cheaply}
The spread is heavily concentrated in the leading stages: stage $0$ accounts
for $80.3\%$ of it and stage $1$ for a further $6.8\%$, while every stage
beyond the second contributes under $0.5\%$. Compensating only the head is
therefore sufficient, and the head gain takes just three distinct values, so
compensation reduces to selecting between three \emph{constants}:
\begin{equation}
  K_{\text{head}}^{-1} \in \{1,\ 0.707107,\ 0.554700\}
  \quad\text{for}\quad |v_0| \in \{0,\ 1,\ 1.5\}.
\end{equation}
Each is realised as a shift-add network evaluated at full width and rounded
once --- $181/256$ and $142/256$, with approximation errors of
$-7.6\times10^{-5}$ and $-1.3\times10^{-5}$ --- requiring four and three adders
respectively, rather than a multiplier. Table~\ref{tab:gaincomp} reports the
resulting accuracy and cost.

\begin{table}[t]
\caption{Gain compensation depth versus residual error and cost.
Perplexity is WikiText-2, full test split, $\mathrm{ctx} = 2048$.}
\label{tab:gaincomp}
\centering
\begin{tabular}{lccrr}
\toprule
Compensated & Residual & \multicolumn{2}{c}{$\Delta$ perplexity} & Area \\
stages & spread & LLaMA-2-7B & Qwen2-7B & ($\mu$m$^2$) \\
\midrule
none      & $93.4\%$ & $+908.3\%$ & $+76.2\%$ & $0$ \\
0 (H1)    & $7.3\%$  & $+0.62\%$  & $+0.77\%$ & <<TBD>> \\
0--1 (H2) & $0.47\%$ & $+0.005\%$ & $+0.008\%$ & <<TBD>> \\
\bottomrule
\end{tabular}
\end{table}

Compensating the leading stage alone holds the end-to-end penalty below
$0.8\%$ on both models, and extending to two stages makes it negligible. The
practical consequence is that the merged datapath remains viable, but its area
advantage must be quoted net of this block --- a distinction the comparison in
Section~\ref{sec:results} makes explicit.
